%2multibyte Version: 5.50.0.2960 CodePage: 65001
% Master Plan del corso "Metodi Quantitativi per la Finanza"
% Versione predisposta per Scientific Workplace 5.5.
% Scelte LaTeX conservative: pacchetti standard, tabelle semplici,
% accenti italiani in forma esplicita, ad esempio probabilit\`{a}.

\documentclass[12pt,a4paper]{article}%
\usepackage{amsmath}
\usepackage{amssymb}
\usepackage{array}
\usepackage{longtable}
\usepackage{amsfonts}
\usepackage{minted}
\usepackage{graphicx}%
\setcounter{MaxMatrixCols}{30}
%TCIDATA{OutputFilter=latex2.dll}
%TCIDATA{Version=5.50.0.2960}
%TCIDATA{Codepage=65001}
%TCIDATA{LastRevised=Tuesday, May 12, 2026 09:39:01}
%TCIDATA{<META NAME="GraphicsSave" CONTENT="32">}
%TCIDATA{<META NAME="SaveForMode" CONTENT="1">}
%TCIDATA{BibliographyScheme=Manual}
%BeginMSIPreambleData
\providecommand{\U}[1]{\protect\rule{.1in}{.1in}}
%EndMSIPreambleData
\usemintedstyle{trac}
\setlength{\textwidth}{15.5cm}
\setlength{\textheight}{23.0cm}
\setlength{\oddsidemargin}{0.3cm}
\setlength{\evensidemargin}{0.3cm}
\setlength{\topmargin}{-0.4cm}
\setlength{\parindent}{0pt}
\setlength{\parskip}{6pt}
\newcommand{\E}{\mathbb{E}}
\newcommand{\Pprob}{\mathbb{P}}
\newcommand{\Qprob}{\mathbb{Q}}
\newcommand{\R}{\mathbb{R}}
\newcommand{\Var}{\operatorname{Var}}
\newcommand{\Cov}{\operatorname{Cov}}
\newcommand{\VaR}{\operatorname{VaR}}
\newcommand{\CVaR}{\operatorname{CVaR}}
\newcommand{\EVPI}{\operatorname{EVPI}}
\newcommand{\VSS}{\operatorname{VSS}}
\begin{document}

\tableofcontents

\begin{center}
{\Large \textbf{Master Plan del corso}}\\[6pt]{\large \textbf{Metodi
Quantitativi per la Finanza}}\\[12pt]
\end{center}

\textbf{Destinatari:} studenti del quinto anno del corso di laurea in Banca e
Risk Management. Principi di Python.\newline\textbf{Lingua del materiale:}
italiano.\newline\textbf{Formato di scrittura:} LaTeX compatibile con
Scientific Workplace 5.5.\newline\textbf{Prodotti finali:} manuale del corso;
slides delle 16 lezioni; esercizi teorici; applicazioni Python; grafici didattici.

\newpage

\section{Impostazione generale del progetto}

\inputminted{python}{Pippo.py}

Il corso deve essere progettato come un percorso coerente su modellizzazione
dell'incertezza, valutazione finanziaria e decisione quantitativa sotto
rischio. La sequenza delle lezioni deve evitare l'impressione di una raccolta
eterogenea di strumenti matematici e deve invece costruire progressivamente un
linguaggio comune.

La progressione concettuale proposta \`{e} la seguente:

\begin{enumerate}
\item rappresentazione probabilistica dell'incertezza;

\item variabili casuali, distribuzioni e momenti;

\item informazione e valore atteso condizionato;

\item processi stocastici e scenari;

\item catene di Markov e rischio di credito;

\item misure di rischio;

\item alberi binomiali e pricing finanziario;

\item programmazione lineare e dualit\`{a};

\item CVaR come problema di ottimizzazione lineare;

\item programmazione stocastica e decisioni in presenza di scenari.
\end{enumerate}

Il livello didattico desiderato \`{e}: rigoroso nella notazione, selettivo
nelle dimostrazioni, applicativo nell'interpretazione finanziaria. Gli
studenti devono poter seguire il passaggio dalla formulazione matematica al
significato economico-finanziario, con particolare attenzione alla coerenza
dei simboli, alla precisione delle ipotesi e all'uso dei grafici come supporto
alla comprensione.

\section{Criteri generali di scrittura}

\subsection{Manuale}

Il manuale deve essere il riferimento scientifico del corso. Ogni capitolo
dovrebbe contenere:

\begin{enumerate}
\item obiettivi della lezione;

\item motivazione finanziaria;

\item definizioni e notazione;

\item risultati teorici principali;

\item esempi numerici;

\item interpretazioni grafiche;

\item applicazioni finanziarie;

\item esercizi svolti;

\item esercizi proposti;

\item sintesi finale.
\end{enumerate}

Per le lezioni Python si aggiungono:

\begin{enumerate}
\item formulazione matematica del problema computazionale;

\item descrizione dell'algoritmo;

\item codice Python commentato;

\item output numerici;

\item grafici didattici;

\item interpretazione economico-finanziaria dei risultati;

\item esercizi di modifica del codice.
\end{enumerate}

\subsection{Slides}

Le slides devono essere uno strumento di lezione, non una versione compressa
del manuale. Per ogni lezione si prevede orientativamente:

\begin{enumerate}
\item motivazione e obiettivi;

\item concetti teorici essenziali;

\item esempi numerici guidati;

\item grafici esplicativi;

\item esercizi da svolgere in aula;

\item sintesi conclusiva.
\end{enumerate}

Per le lezioni Python le slides devono privilegiare formulazione, algoritmo,
codice selezionato, output e interpretazione.

\section{Convenzioni preliminari di notazione}

Le convenzioni definitive saranno raccolte nel documento separato
\texttt{MQF\_Notazione.tex}. In prima approssimazione:

\begin{center}%
\begin{tabular}
[c]{cp{8.5cm}}\hline
\textbf{Simbolo} & \textbf{Significato}\\\hline
\multicolumn{1}{l}{$\left(  \Omega,\mathcal{F},\mathbb{P}\right)  $} & Spazio
di probabilit\`{a}.\\
\multicolumn{1}{l}{$X,Y,Z$} & Variabili casuali.\\
\multicolumn{1}{l}{$\mathbb{E}[X]$, $\operatorname{Var}(X)$,
$\operatorname{Cov}(X,Y)$} & Valore atteso, varianza e covarianza.\\
\multicolumn{1}{l}{$\{\mathcal{F}_{t}\}_{t=0}^{T}$} & Filtrazione in tempo
discreto.\\
\multicolumn{1}{l}{$S_{t}$} & Prezzo di un'attivit\`{a} finanziaria al tempo
$t$.\\
\multicolumn{1}{l}{$r$} & Tasso privo di rischio su un periodo, salvo diversa
specificazione.\\
\multicolumn{1}{l}{$\mathbb{Q}$} & Misura risk-neutral, quando introdotta.\\
\multicolumn{1}{l}{$x$} & Vettore di variabili decisionali.\\
\multicolumn{1}{l}{$\xi$} & Vettore aleatorio oppure informazione incerta nei
problemi decisionali.\\
\multicolumn{1}{l}{$\omega_{s}$, $p_{s}$} & Scenario $s$ e probabilit\`{a}
dello scenario $s$.\\\hline
\end{tabular}

\end{center}

Salvo diversa indicazione, il corso lavora prevalentemente in tempo discreto:
\[
t=0,1,\ldots,T.
\]


\section{Struttura delle 16 lezioni}

\begin{longtable}{>{\centering}p{0.9cm}>{\centering}p{0.8cm}p{4.3cm}p{4.2cm}p{3.3cm}}
\hline
\textbf{Lez.} & \textbf{Tipo} & \textbf{Titolo definitivo} & \textbf{Tema centrale} & \textbf{Output principale} \\
\hline
\endfirsthead
\hline
\textbf{Lez.} & \textbf{Tipo} & \textbf{Titolo definitivo} & \textbf{Tema centrale} & \textbf{Output principale} \\
\hline
\endhead
1 & P & Elementi di teoria della probabilit\`{a} & Spazio di probabilit\`{a}, eventi, probabilit\`{a} condizionata & Capitolo teorico ed esercizi \\
2 & P & Variabili casuali e distribuzioni & Variabili casuali, distribuzioni, momenti & Capitolo teorico, esercizi e grafici \\
3 & P & Valori attesi condizionati e informazione & Condizionamento, partizioni, informazione & Capitolo teorico ed esercizi \\
4 & P & Processi stocastici, traiettorie e scenari & Processi in tempo discreto, simulazione, scenari & Capitolo teorico e grafici \\
5 & C & Applicazioni Python a processi stocastici e valori attesi condizionati & Simulazione e calcolo computazionale & Notebook o script e slides applicative \\
6 & P & Catene di Markov: definizioni e propriet\`{a} & Stati, transizioni, matrice di transizione & Capitolo teorico ed esercizi \\
7 & P & Catene di Markov e misure di rischio & Evoluzione delle distribuzioni, VaR, CVaR & Capitolo teorico, esercizi e grafici \\
8 & C & Applicazioni Python al rischio di credito & Transizioni di rating e rischio di default & Notebook o script e slides applicative \\
9 & P & Variabili casuali binomiali e alberi binomiali & Modello binomiale, payoff, probabilit\`{a} risk-neutral & Capitolo teorico ed esercizi \\
10 & C & Applicazioni Python al pricing di opzioni e obbligazioni & Pricing numerico in alberi discreti & Notebook o script e slides applicative \\
11 & P & Programmazione lineare & Formulazione, vincoli, regione ammissibile & Capitolo teorico, esercizi e grafici \\
12 & P & Dualit\`{a} nella programmazione lineare & Problema duale, prezzi ombra, interpretazione economica & Capitolo teorico ed esercizi \\
13 & C & Applicazioni Python di programmazione lineare: calcolo del CVaR & CVaR come problema di programmazione lineare & Notebook o script e slides applicative \\
14 & P & Programmazione lineare stocastica I & Decisioni here-and-now, scenari, recourse & Capitolo teorico ed esercizi \\
15 & P & Programmazione lineare stocastica II & Valore dell'informazione, decomposizione concettuale & Capitolo teorico ed esercizi \\
16 & C & Applicazioni Python di programmazione stocastica & Implementazione di modelli stocastici discreti & Notebook o script e slides applicative \\
\hline
\end{longtable}


\textbf{Legenda:} P = lezione prevalentemente teorica; C = lezione con forte
componente computazionale o applicativa.

\section{Schede progettuali delle lezioni}

\subsection{Lezione 1 -- Elementi di teoria della probabilit\`{a}}

\textbf{Tipo:} P.\newline\textbf{Domanda guida:} come si rappresenta
formalmente l'incertezza di un payoff finanziario?\newline%
\textbf{Prerequisiti:} insiemi, algebra elementare degli eventi, nozioni di
probabilit\`{a} di base.\newline\textbf{Obiettivi didattici:} introdurre lo
spazio di probabilit\`{a}; distinguere eventi, probabilit\`{a} marginali,
congiunte e condizionate; collegare il formalismo probabilistico alla nozione
finanziaria di rischio.\newline\textbf{Notazione introdotta:} $\Omega$,
$\mathcal{F}$, $\mathbb{P}$, eventi $A,B$, probabilit\`{a} condizionata
$\mathbb{P}(A\mid B)$.\newline\textbf{Formule principali:}
\[
\mathbb{P}(A\mid B)=\frac{\mathbb{P}(A\cap B)}{\mathbb{P}(B)}, \qquad
\mathbb{P}(A)=\sum_{i} \mathbb{P}(A\mid B_{i})\mathbb{P}(B_{i}),
\]
\[
\mathbb{P}(B_{j}\mid A)=\frac{\mathbb{P}(A\mid B_{j})\mathbb{P}(B_{j})}%
{\sum_{i} \mathbb{P}(A\mid B_{i})\mathbb{P}(B_{i})}.
\]
\textbf{Esercizi previsti:} probabilit\`{a} condizionata su eventi di mercato;
aggiornamento bayesiano di uno scenario di default; calcolo di probabilit\`{a}
congiunte.\newline\textbf{Grafici previsti:} diagrammi di Venn; albero
probabilistico a due stadi.\newline\textbf{Applicazioni Python collegate:}
Lezione 5.\newline\textbf{Stato di avanzamento:} bozza progettuale
predisposta; contenuti teorici da sviluppare.

\subsection{Lezione 2 -- Variabili casuali e distribuzioni}

\textbf{Tipo:} P.\newline\textbf{Domanda guida:} come si descrive
quantitativamente un rendimento, una perdita o un payoff aleatorio?\newline%
\textbf{Prerequisiti:} Lezione 1.\newline\textbf{Obiettivi didattici:}
definire variabili casuali discrete e continue; introdurre distribuzioni,
funzione di ripartizione, momenti e quantili; collegare momenti e
rischio.\newline\textbf{Notazione introdotta:} $X$, $F_{X}(x)$, $f_{X}(x)$,
$p_{X}(x)$, $\mathbb{E}[X]$, $\operatorname{Var}(X)$.\newline\textbf{Formule
principali:}
\[
F_{X}(x)=\mathbb{P}(X\leq x), \qquad\mathbb{E}[X]=\sum_{x} x p_{X}(x),
\qquad\operatorname{Var}(X)=\mathbb{E}\left[  (X-\mathbb{E}[X])^{2}\right]  .
\]
\textbf{Esercizi previsti:} payoff discreto; rendimento atteso e varianza;
quantile di una distribuzione discreta di perdite.\newline\textbf{Grafici
previsti:} funzione di massa; densit\`{a}; funzione di ripartizione;
rappresentazione dei quantili.\newline\textbf{Applicazioni Python collegate:}
Lezioni 5, 8 e 13.\newline\textbf{Stato di avanzamento:} bozza progettuale
predisposta; esempi numerici da selezionare.

\subsection{Lezione 3 -- Valori attesi condizionati e informazione}

\textbf{Tipo:} P.\newline\textbf{Domanda guida:} come cambia la valutazione di
un payoff quando arriva nuova informazione?\newline\textbf{Prerequisiti:}
Lezioni 1--2.\newline\textbf{Obiettivi didattici:} introdurre il valore atteso
condizionato rispetto a eventi, partizioni e sigma-algebre; interpretarlo come
migliore previsione data l'informazione disponibile.\newline\textbf{Notazione
introdotta:} $\mathbb{E}[X\mid A]$, $\mathbb{E}[X\mid\mathcal{G}]$, partizioni
finite, sigma-algebre informative.\newline\textbf{Formule principali:}
\[
\mathbb{E}[X\mid A]=\sum_{x} x\,\mathbb{P}(X=x\mid A), \qquad\mathbb{E}\left[
\mathbb{E}[X\mid\mathcal{G}]\right]  =\mathbb{E}[X],
\]
\[
\mathbb{E}\left[  \mathbb{E}[X\mid\mathcal{G}_{2}]\mid\mathcal{G}_{1}\right]
=\mathbb{E}[X\mid\mathcal{G}_{1}], \qquad\mathcal{G}_{1}\subseteq
\mathcal{G}_{2}.
\]
\textbf{Esercizi previsti:} calcolo su partizioni; aggiornamento di payoff
condizionati a stati di mercato; applicazione della propriet\`{a} della
torre.\newline\textbf{Grafici previsti:} albero informativo; rappresentazione
per stati e partizioni.\newline\textbf{Applicazioni Python collegate:} Lezione
5; Lezione 16 per la distinzione tra informazione anticipativa e non
anticipativa.\newline\textbf{Stato di avanzamento:} bozza progettuale
predisposta; dimostrazioni essenziali da definire.

\subsection{Lezione 4 -- Processi stocastici, traiettorie e scenari}

\textbf{Tipo:} P.\newline\textbf{Domanda guida:} come si rappresenta
l'evoluzione temporale di prezzi, rendimenti o fattori di rischio?\newline%
\textbf{Prerequisiti:} Lezioni 1--3.\newline\textbf{Obiettivi didattici:}
introdurre processi stocastici in tempo discreto; distinguere traiettorie,
stati e scenari; collegare simulazione e rappresentazione
dell'incertezza.\newline\textbf{Notazione introdotta:} $\{X_{t}\}_{t=0}^{T}$,
$\{S_{t}\}_{t=0}^{T}$, $\mathcal{F}_{t}$, scenari $\omega_{s}$.\newline%
\textbf{Formule principali:}
\[
X_{t}:\Omega\rightarrow\mathbb{R}, \qquad\{X_{t}\}_{t=0}^{T}, \qquad X_{t} \;
\hbox{\`{e} } \mathcal{F}_{t}\hbox{-misurabile}.
\]
\textbf{Esercizi previsti:} costruzione di traiettorie discrete; calcolo di
medie e varianze per stati temporali; identificazione dell'informazione
disponibile.\newline\textbf{Grafici previsti:} traiettorie simulate; ventaglio
di scenari; albero degli scenari.\newline\textbf{Applicazioni Python
collegate:} Lezione 5; Lezioni 8, 10 e 16.\newline\textbf{Stato di
avanzamento:} bozza progettuale predisposta; scelta dei processi
esemplificativi da completare.

\subsection{Lezione 5 -- Applicazioni Python a processi stocastici e valori
attesi condizionati}

\textbf{Tipo:} C.\newline\textbf{Domanda guida:} come si simulano traiettorie
e come si stimano quantitativamente valori attesi condizionati?\newline%
\textbf{Prerequisiti:} Lezioni 1--4; basi operative di Python.\newline%
\textbf{Obiettivi didattici:} implementare simulazioni di processi discreti;
stimare valori attesi condizionati; visualizzare traiettorie e
scenari.\newline\textbf{Notazione introdotta:} nessuna notazione teorica
sostanziale; raccordo tra simboli matematici e variabili Python.\newline%
\textbf{Formule principali:}
\[
\widehat{\mathbb{E}}[X]=\frac{1}{N}\sum_{n=1}^{N}X^{(n)}, \qquad
\widehat{\mathbb{E}}[X\mid G=g]=\frac{1}{N_{g}}\sum_{n:G^{(n)}=g}X^{(n)}.
\]
\textbf{Esercizi previsti:} modificare parametri di simulazione; aumentare il
numero di scenari; interpretare la stabilit\`{a} numerica.\newline%
\textbf{Grafici previsti:} traiettorie multiple; media temporale; bande
empiriche.\newline\textbf{Applicazioni Python collegate:} applicazione Python
1.\newline\textbf{Stato di avanzamento:} bozza progettuale predisposta; codice
da scrivere.

\subsection{Lezione 6 -- Catene di Markov: definizioni e propriet\`{a}}

\textbf{Tipo:} P.\newline\textbf{Domanda guida:} come si modellano transizioni
probabilistiche tra stati finanziari, ad esempio classi di rating?\newline%
\textbf{Prerequisiti:} Lezioni 1--4.\newline\textbf{Obiettivi didattici:}
definire catene di Markov in tempo discreto; introdurre matrice di
transizione; analizzare distribuzioni a uno o pi\`{u} passi.\newline%
\textbf{Notazione introdotta:} stati $i,j$, matrice $P=(p_{ij})$,
distribuzione iniziale $\pi_{0}$, distribuzione $\pi_{t}$.\newline%
\textbf{Formule principali:}
\[
p_{ij}=\mathbb{P}(X_{t+1}=j\mid X_{t}=i), \qquad\mathbb{P}(X_{t+1}=j\mid
X_{t}=i,X_{t-1},\ldots,X_{0})=p_{ij},
\]
\[
\pi_{t}=\pi_{0} P^{t}.
\]
\textbf{Esercizi previsti:} calcolo di transizioni a due periodi;
distribuzione futura di rating; probabilit\`{a} di default assorbente.\newline%
\textbf{Grafici previsti:} grafo degli stati; heatmap della matrice di
transizione; evoluzione delle probabilit\`{a} di stato.\newline%
\textbf{Applicazioni Python collegate:} Lezione 8.\newline\textbf{Stato di
avanzamento:} bozza progettuale predisposta; matrice di transizione
esemplificativa da definire.

\subsection{Lezione 7 -- Catene di Markov e misure di rischio}

\textbf{Tipo:} P.\newline\textbf{Domanda guida:} come si passa dalla dinamica
degli stati alla misurazione del rischio di perdita?\newline%
\textbf{Prerequisiti:} Lezioni 2 e 6.\newline\textbf{Obiettivi didattici:}
collegare distribuzioni di perdita, VaR e CVaR; usare catene di Markov per
generare distribuzioni future di rischio.\newline\textbf{Notazione
introdotta:} perdita $L$, $\operatorname{VaR}_{\alpha}(L)$,
$\operatorname{CVaR}_{\alpha}(L)$.\newline\textbf{Formule principali:}
\[
\operatorname{VaR}_{\alpha}(L)=\inf\{\ell:\mathbb{P}(L\leq\ell)\geq\alpha\},
\]
\[
\operatorname{CVaR}_{\alpha}(L)=\mathbb{E}[L\mid L\geq\operatorname{VaR}%
_{\alpha}(L)]
\]
nel caso discreto elementare in cui la coda sia trattata senza correzioni di
interpolazione.\newline\textbf{Esercizi previsti:} calcolo di VaR e CVaR su
perdite discrete; confronto tra due portafogli; interpretazione delle
code.\newline\textbf{Grafici previsti:} distribuzione delle perdite; quantile
VaR; area di coda CVaR.\newline\textbf{Applicazioni Python collegate:} Lezioni
8 e 13.\newline\textbf{Stato di avanzamento:} bozza progettuale predisposta;
precisare definizione operativa del CVaR discreto.

\subsection{Lezione 8 -- Applicazioni Python al rischio di credito}

\textbf{Tipo:} C.\newline\textbf{Domanda guida:} come si implementa un modello
di transizione di rating e si stima il rischio di credito?\newline%
\textbf{Prerequisiti:} Lezioni 6--7; basi operative di Python.\newline%
\textbf{Obiettivi didattici:} implementare una matrice di transizione;
simulare evoluzioni di rating; calcolare probabilit\`{a} di default,
distribuzioni di perdita e misure di rischio.\newline\textbf{Notazione
introdotta:} raccordo computazionale con $P$, $\pi_{t}$, $L$,
$\operatorname{VaR}_{\alpha}(L)$ e $\operatorname{CVaR}_{\alpha}(L)$%
.\newline\textbf{Formule principali:}
\[
\pi_{t}=\pi_{0}P^{t}, \qquad\widehat{\mathbb{P}}(\hbox{default entro }T)=\frac
{1}{N}\sum_{n=1}^{N} I_{n},
\]
dove $I_{n}$ indica l'evento di default simulato nella traiettoria
$n$.\newline\textbf{Esercizi previsti:} modificare la matrice di transizione;
introdurre uno stato assorbente; confrontare portafogli con diverse
esposizioni.\newline\textbf{Grafici previsti:} heatmap della matrice;
probabilit\`{a} di default nel tempo; distribuzione simulata delle
perdite.\newline\textbf{Applicazioni Python collegate:} applicazione Python
2.\newline\textbf{Stato di avanzamento:} bozza progettuale predisposta; codice
da scrivere.

\subsection{Lezione 9 -- Variabili casuali binomiali e alberi binomiali}

\textbf{Tipo:} P.\newline\textbf{Domanda guida:} come si discretizza
l'evoluzione di un prezzo finanziario per valutare un derivato?\newline%
\textbf{Prerequisiti:} Lezioni 1--4.\newline\textbf{Obiettivi didattici:}
introdurre variabili binomiali, alberi dei prezzi, payoff terminali e
probabilit\`{a} risk-neutral.\newline\textbf{Notazione introdotta:} $u$, $d$,
$q$, $S_{t}$, $V_{t}$, payoff $H$.\newline\textbf{Formule principali:}
\[
S_{t+1}=%
\begin{cases}
uS_{t}, & \hbox{stato up},\\
dS_{t}, & \hbox{stato down},
\end{cases}
\qquad q=\frac{(1+r)-d}{u-d},
\]
\[
V_{t}=\frac{1}{1+r}\left(  qV_{t+1}^{u}+(1-q)V_{t+1}^{d}\right)  .
\]
\textbf{Esercizi previsti:} pricing di opzione call a un periodo; costruzione
di albero a due periodi; calcolo di payoff terminali e valore
iniziale.\newline\textbf{Grafici previsti:} albero binomiale dei prezzi;
albero dei payoff; backward induction visuale.\newline\textbf{Applicazioni
Python collegate:} Lezione 10.\newline\textbf{Stato di avanzamento:} bozza
progettuale predisposta; ipotesi di assenza di arbitraggio da formalizzare.

\subsection{Lezione 10 -- Applicazioni Python al pricing di opzioni e
obbligazioni}

\textbf{Tipo:} C.\newline\textbf{Domanda guida:} come si implementa
numericamente il pricing ricorsivo in un albero discreto?\newline%
\textbf{Prerequisiti:} Lezione 9; basi operative di Python.\newline%
\textbf{Obiettivi didattici:} costruire alberi binomiali in Python; calcolare
payoff; implementare backward induction; estendere il metodo a strumenti
obbligazionari semplici.\newline\textbf{Notazione introdotta:} raccordo
computazionale con $S_{t}$, $V_{t}$, $u$, $d$, $q$, $r$.\newline%
\textbf{Formule principali:}
\[
H=(S_{T}-K)^{+}, \qquad V_{t}=\frac{1}{1+r}\mathbb{E}^{\mathbb{Q}}[V_{t+1}%
\mid\mathcal{F}_{t}].
\]
\textbf{Esercizi previsti:} modificare strike, tasso e volatilit\`{a};
confrontare prezzi; introdurre un payoff alternativo.\newline\textbf{Grafici
previsti:} albero del prezzo; albero del valore; sensibilit\`{a} del prezzo
rispetto ai parametri.\newline\textbf{Applicazioni Python collegate:}
applicazione Python 3.\newline\textbf{Stato di avanzamento:} bozza progettuale
predisposta; codice da scrivere.

\subsection{Lezione 11 -- Programmazione lineare}

\textbf{Tipo:} P.\newline\textbf{Domanda guida:} come si formula e risolve una
decisione finanziaria vincolata?\newline\textbf{Prerequisiti:} algebra lineare
elementare; ottimizzazione di base.\newline\textbf{Obiettivi didattici:}
introdurre variabili decisionali, funzione obiettivo, vincoli, forma standard,
regione ammissibile e soluzione ottima.\newline\textbf{Notazione introdotta:}
$x$, $c$, $A$, $b$, problema primale.\newline\textbf{Formule principali:}
\[
\max_{x} \; c^{\prime}x \qquad\hbox{soggetto a} \qquad Ax\leq b, \quad
x\geq0.
\]
In forma alternativa:
\[
\min_{x} \; c^{\prime}x \qquad\hbox{soggetto a} \qquad Ax=b, \quad x\geq0.
\]
\textbf{Esercizi previsti:} problema di allocazione di portafoglio con
vincoli; rappresentazione grafica in due variabili; identificazione della
soluzione ottima.\newline\textbf{Grafici previsti:} regione ammissibile; rette
di livello; soluzione ottima su un vertice.\newline\textbf{Applicazioni Python
collegate:} Lezione 13 e Lezione 16.\newline\textbf{Stato di avanzamento:}
bozza progettuale predisposta; scegliere convenzione primaria tra forma di
massimo e forma di minimo.

\subsection{Lezione 12 -- Dualit\`{a} nella programmazione lineare}

\textbf{Tipo:} P.\newline\textbf{Domanda guida:} qual \`{e} il significato
economico dei prezzi ombra associati ai vincoli?\newline\textbf{Prerequisiti:}
Lezione 11.\newline\textbf{Obiettivi didattici:} introdurre il problema duale;
interpretare variabili duali e prezzi ombra; collegare primalit\`{a},
dualit\`{a} e sensibilit\`{a}.\newline\textbf{Notazione introdotta:} variabili
duali $y$, problema primale, problema duale, valore ottimo $z^{*}$%
.\newline\textbf{Formule principali:}
\[
\max_{x} \; c^{\prime}x \quad\hbox{soggetto a} \quad Ax\leq b,\; x\geq0,
\]
\[
\min_{y} \; b^{\prime}y \quad\hbox{soggetto a} \quad A^{\prime}y\geq c,\;
y\geq0.
\]
\textbf{Esercizi previsti:} derivazione del duale da un problema semplice;
interpretazione dei moltiplicatori; analisi di variazioni nel vincolo.\newline%
\textbf{Grafici previsti:} interpretazione geometrica primal-dual; rette di
supporto; valore marginale dei vincoli.\newline\textbf{Applicazioni Python
collegate:} Lezioni 13 e 16.\newline\textbf{Stato di avanzamento:} bozza
progettuale predisposta; esempi economico-finanziari da scegliere.

\subsection{Lezione 13 -- Applicazioni Python di programmazione lineare:
calcolo del CVaR}

\textbf{Tipo:} C.\newline\textbf{Domanda guida:} come si calcola il CVaR di
portafoglio attraverso un problema di programmazione lineare?\newline%
\textbf{Prerequisiti:} Lezioni 7, 11 e 12; basi operative di Python.\newline%
\textbf{Obiettivi didattici:} formulare il CVaR come problema di
programmazione lineare; implementare il modello; interpretare la soluzione e i
vincoli.\newline\textbf{Notazione introdotta:} livello $\alpha$, soglia $\eta
$, variabili ausiliarie $z_{s}$, scenari $s$.\newline\textbf{Formule
principali:}
\[
\min_{x,\eta,z} \; \eta+\frac{1}{1-\alpha}\sum_{s=1}^{S}p_{s} z_{s}
\]
\[
\hbox{soggetto a} \qquad z_{s}\geq L_{s}(x)-\eta, \quad z_{s}\geq0, \quad
s=1,\ldots,S,
\]
con eventuali vincoli di portafoglio su $x$.\newline\textbf{Esercizi
previsti:} modificare il livello di confidenza; confrontare portafogli;
interpretare le variabili ausiliarie.\newline\textbf{Grafici previsti:}
distribuzione delle perdite; VaR e CVaR; sensibilit\`{a} rispetto ad $\alpha
$.\newline\textbf{Applicazioni Python collegate:} applicazione Python
4.\newline\textbf{Stato di avanzamento:} bozza progettuale predisposta;
libreria di ottimizzazione Python da decidere.

\subsection{Lezione 14 -- Programmazione lineare stocastica I}

\textbf{Tipo:} P.\newline\textbf{Domanda guida:} come si prende una decisione
oggi quando gli esiti futuri sono incerti?\newline\textbf{Prerequisiti:}
Lezioni 3, 4, 11 e 12.\newline\textbf{Obiettivi didattici:} introdurre
scenari, decisioni here-and-now, decisioni wait-and-see e recourse; formulare
problemi stocastici discreti a due stadi.\newline\textbf{Notazione
introdotta:} scenari $s$, probabilit\`{a} $p_{s}$, decisione di primo stadio
$x$, decisione di secondo stadio $y_{s}$.\newline\textbf{Formule principali:}
\[
\min_{x} \; c^{\prime}x+\mathbb{E}_{\xi}[Q(x,\xi)]
\]
con approssimazione discreta
\[
\min_{x,y_{s}} \; c^{\prime}x+\sum_{s=1}^{S}p_{s} q_{s}^{\prime}y_{s}
\]
\[
\hbox{soggetto a} \qquad Ax=b, \quad T_{s} x+W_{s} y_{s}=h_{s}, \quad x\geq0,
\quad y_{s}\geq0.
\]
\textbf{Esercizi previsti:} formulazione di un problema a due stadi;
costruzione di tabella scenari; distinzione tra decisioni anticipative e non
anticipative.\newline\textbf{Grafici previsti:} albero degli scenari; schema
here-and-now/recourse; confronto tra scenari.\newline\textbf{Applicazioni
Python collegate:} Lezione 16.\newline\textbf{Stato di avanzamento:} bozza
progettuale predisposta; esempio finanziario di riferimento da scegliere.

\subsection{Lezione 15 -- Programmazione lineare stocastica II}

\textbf{Tipo:} P.\newline\textbf{Domanda guida:} qual \`{e} il valore
economico dell'informazione e della modellizzazione stocastica?\newline%
\textbf{Prerequisiti:} Lezione 14.\newline\textbf{Obiettivi didattici:}
approfondire valore dell'informazione, soluzione wait-and-see, expected value
solution, recourse problem e non anticipativit\`{a}.\newline\textbf{Notazione
introdotta:} $z^{SP}$, $z^{WS}$, $z^{EV}$, $\operatorname{EVPI}$,
$\operatorname{VSS}$.\newline\textbf{Formule principali:}
\[
z^{SP}=\hbox{valore ottimo del problema stocastico}, \qquad z^{WS}=\sum
_{s=1}^{S}p_{s} z_{s}^{WS},
\]
\[
\operatorname{EVPI}=z^{SP}-z^{WS}, \qquad\operatorname{VSS}=z^{EV}-z^{SP},
\]
per una convenzione di minimizzazione.\newline\textbf{Esercizi previsti:}
calcolo di EVPI in un esempio discreto; confronto tra soluzione deterministica
equivalente e soluzione stocastica; interpretazione economica.\newline%
\textbf{Grafici previsti:} confronto dei valori ottimi; albero con decisioni
non anticipative; diagramma dei livelli informativi.\newline%
\textbf{Applicazioni Python collegate:} Lezione 16.\newline\textbf{Stato di
avanzamento:} bozza progettuale predisposta; convenzioni di segno per EVPI e
VSS da confermare.

\subsection{Lezione 16 -- Applicazioni Python di programmazione stocastica}

\textbf{Tipo:} C.\newline\textbf{Domanda guida:} come si implementa un modello
stocastico discreto e come si interpretano le decisioni ottime?\newline%
\textbf{Prerequisiti:} Lezioni 14--15; basi operative di Python.\newline%
\textbf{Obiettivi didattici:} costruire un modello stocastico a scenari;
risolverlo numericamente; confrontare soluzione stocastica, soluzione
deterministica equivalente e wait-and-see.\newline\textbf{Notazione
introdotta:} raccordo computazionale con $x$, $y_{s}$, $p_{s}$, $z^{SP}$,
$z^{WS}$, $z^{EV}$, $\operatorname{EVPI}$ e $\operatorname{VSS}$%
.\newline\textbf{Formule principali:}
\[
\min_{x,y_{s}} \; c^{\prime}x+\sum_{s=1}^{S}p_{s} q_{s}^{\prime}y_{s},
\]
con vincoli di primo stadio, vincoli di recourse e vincoli di non
anticipativit\`{a}.\newline\textbf{Esercizi previsti:} modificare
probabilit\`{a} degli scenari; introdurre vincoli di budget o rischio;
confrontare decisioni ottime.\newline\textbf{Grafici previsti:} albero degli
scenari; confronto tra soluzioni; distribuzione dei risultati per
scenario.\newline\textbf{Applicazioni Python collegate:} applicazione Python
5.\newline\textbf{Stato di avanzamento:} bozza progettuale predisposta; codice
da scrivere.

\section{Registro dei grafici didattici previsti}

\begin{longtable}{p{4.0cm}p{4.2cm}p{6.1cm}}
\hline
\textbf{Tema} & \textbf{Grafico} & \textbf{Uso didattico} \\
\hline
\endfirsthead
\hline
\textbf{Tema} & \textbf{Grafico} & \textbf{Uso didattico} \\
\hline
\endhead
Probabilit\`{a} condizionata & Albero probabilistico & Visualizzare eventi sequenziali e aggiornamento informativo. \\
Variabili casuali & Funzione di massa o densit\`{a} & Collegare distribuzione, media e dispersione. \\
Quantili & Funzione di ripartizione & Introdurre VaR e code della distribuzione. \\
Valore atteso condizionato & Albero informativo & Mostrare il ruolo dell'informazione. \\
Processi stocastici & Traiettorie simulate & Visualizzare dinamica e incertezza. \\
Scenari & Ventaglio di scenari & Collegare simulazione e decisione. \\
Catene di Markov & Grafo degli stati & Visualizzare transizioni. \\
Matrici di transizione & Heatmap & Evidenziare persistenza e probabilit\`{a} di default. \\
Misure di rischio & Distribuzione delle perdite & Confrontare VaR e CVaR. \\
Alberi binomiali & Albero dei prezzi e dei payoff & Spiegare pricing ricorsivo. \\
Programmazione lineare & Regione ammissibile & Interpretare vincoli e ottimo. \\
Dualit\`{a} & Rette di supporto e prezzi ombra & Interpretare valori marginali. \\
CVaR via PL & Perdita e variabili ausiliarie & Collegare rischio e ottimizzazione. \\
Programmazione stocastica & Albero degli scenari & Distinguere decisioni e informazione. \\
Valore dell'informazione & Confronto SP/WS/EV & Interpretare EVPI e VSS. \\
\hline
\end{longtable}


\section{Registro degli esercizi in aula}

Ogni lezione teorica dovrebbe includere almeno due esercizi in aula: un
esercizio di calcolo o derivazione e un esercizio di interpretazione
finanziaria. Le lezioni Python dovrebbero includere almeno un esercizio di
modifica del codice e un esercizio di interpretazione dell'output.

\begin{longtable}{>{\centering}p{1.1cm}p{4.1cm}p{4.1cm}p{4.5cm}}
\hline
\textbf{Lez.} & \textbf{Esercizio 1} & \textbf{Esercizio 2} & \textbf{Esercizio grafico/computazionale} \\
\hline
\endfirsthead
\hline
\textbf{Lez.} & \textbf{Esercizio 1} & \textbf{Esercizio 2} & \textbf{Esercizio grafico/computazionale} \\
\hline
\endhead
1 & Probabilit\`{a} condizionata & Formula di Bayes & Albero probabilistico \\
2 & Media e varianza di payoff discreto & Quantile di perdita & Grafico della funzione di ripartizione \\
3 & Valore atteso condizionato su partizione & Propriet\`{a} della torre & Albero informativo \\
4 & Traiettorie discrete & Scenari e filtrazione & Ventaglio di scenari \\
5 & Modifica della simulazione & Stima condizionata & Grafico delle traiettorie \\
6 & Transizioni a due periodi & Stato assorbente & Grafo di Markov \\
7 & VaR discreto & CVaR discreto & Distribuzione delle perdite \\
8 & Transizioni di rating & Probabilit\`{a} di default & Heatmap delle transizioni \\
9 & Pricing binomiale a un periodo & Backward induction & Albero binomiale \\
10 & Variazione di strike e tasso & Payoff alternativo & Albero dei prezzi e dei valori \\
11 & Formulazione PL & Soluzione grafica & Regione ammissibile \\
12 & Costruzione del duale & Prezzi ombra & Interpretazione geometrica \\
13 & CVaR via PL & Sensibilit\`{a} ad $\alpha$ & Grafico VaR--CVaR \\
14 & Modello a due stadi & Non anticipativit\`{a} & Albero degli scenari \\
15 & EVPI & VSS & Confronto dei valori \\
16 & Modifica degli scenari & Confronto SP/EV/WS & Grafico dei risultati \\
\hline
\end{longtable}


\section{Registro delle applicazioni Python}

\begin{longtable}{>{\centering}p{1.5cm}>{\centering}p{1.4cm}p{4.4cm}p{6.5cm}}
\hline
\textbf{Applic.} & \textbf{Lez.} & \textbf{Tema} & \textbf{Output computazionale} \\
\hline
\endfirsthead
\hline
\textbf{Applic.} & \textbf{Lez.} & \textbf{Tema} & \textbf{Output computazionale} \\
\hline
\endhead
1 & 5 & Processi stocastici e valori attesi condizionati & Simulazioni, stime condizionate, grafici delle traiettorie. \\
2 & 8 & Rischio di credito & Matrici di transizione, probabilit\`{a} di default, VaR e CVaR. \\
3 & 10 & Pricing di opzioni e obbligazioni & Alberi binomiali, backward induction, analisi di sensibilit\`{a}. \\
4 & 13 & Programmazione lineare e CVaR & Formulazione PL, soluzione numerica, confronto VaR--CVaR. \\
5 & 16 & Programmazione stocastica & Scenari, recourse, EVPI, VSS. \\
\hline
\end{longtable}


Criteri comuni per il codice Python:

\begin{enumerate}
\item codice ben commentato;

\item separazione tra dati, parametri, funzioni e output;

\item uso di nomi coerenti con la notazione matematica;

\item grafici leggibili;

\item interpretazione dei risultati dopo ogni blocco computazionale;

\item esercizi di modifica del codice.
\end{enumerate}

\section{File previsti del progetto}

\subsection{Manuale}
\begin{verbatim}
/manuale
  MQF_Manuale_Master.tex
  MQF_Capitolo_01.tex
  MQF_Capitolo_02.tex
  ...
  MQF_Capitolo_16.tex
\end{verbatim}

\subsection{Slides}
\begin{verbatim}
/slides
  MQF_Slides_01.tex
  MQF_Slides_02.tex
  ...
  MQF_Slides_16.tex
\end{verbatim}

\subsection{Python}
\begin{verbatim}
/python
  MQF_Python_01_Processi_Stocastici.py
  MQF_Python_02_Rischio_Credito.py
  MQF_Python_03_Pricing.py
  MQF_Python_04_CVaR_PL.py
  MQF_Python_05_Programmazione_Stocastica.py
\end{verbatim}

\subsection{Documenti di coordinamento}
\begin{verbatim}
/progetto
  MQF_Master_Plan.tex
  MQF_Notazione.tex
  MQF_Registro_Esercizi.tex
  MQF_Registro_Grafici.tex
  MQF_Registro_Decisioni.tex
\end{verbatim}

\section{Stato di avanzamento generale}

\begin{longtable}{p{4.2cm}p{3.0cm}p{7.0cm}}
\hline
\textbf{Area} & \textbf{Stato} & \textbf{Note} \\
\hline
\endfirsthead
\hline
\textbf{Area} & \textbf{Stato} & \textbf{Note} \\
\hline
\endhead
Struttura delle 16 lezioni & Bozza iniziale & Derivata dal syllabus e uniformata nei titoli. \\
Notazione generale & Da consolidare & Preparare documento separato \texttt{MQF\_Notazione.tex}. \\
Manuale & Non iniziato & Scrittura per capitoli separati. \\
Slides & Non iniziate & Da produrre dopo schema dettagliato di ciascuna lezione. \\
Esercizi teorici & Bozza iniziale & Da sviluppare lezione per lezione. \\
Applicazioni Python & Bozza iniziale & Cinque applicazioni previste dal syllabus. \\
Grafici & Bozza iniziale & Registro grafici da validare. \\
Compatibilit\`{a} SWP 5.5 & Da verificare & Evitare pacchetti LaTeX moderni non necessari. \\
\hline
\end{longtable}


\section{Questioni aperte}

\begin{enumerate}
\item Definire il titolo definitivo del manuale.

\item Definire la struttura LaTeX compatibile con Scientific Workplace 5.5.

\item Stabilire se i grafici saranno prodotti prevalentemente in Python e
inclusi come immagini, oppure scritti in LaTeX quando compatibile.

\item Decidere il livello di dettaglio delle dimostrazioni.

\item Preparare il documento \texttt{MQF\_Notazione.tex}.

\item Preparare un template standard per capitoli del manuale.

\item Preparare un template standard per slides.

\item Definire formato e librerie Python ammesse per le applicazioni computazionali.

\item Decidere se includere soluzioni complete degli esercizi nel manuale o in
un fascicolo separato.

\item Definire criteri di valutazione, se il materiale dovr\`{a} essere
coerente con una prova d'esame.
\end{enumerate}

\section{Prossimi passi consigliati}

\begin{enumerate}
\item Validare o modificare i titoli definitivi delle 16 lezioni.

\item Predisporre il documento \texttt{MQF\_Notazione.tex}.

\item Definire il template LaTeX del manuale.

\item Definire il template LaTeX delle slides.

\item Sviluppare la scheda dettagliata della Lezione 1.

\item Solo dopo la validazione della Lezione 1, replicare lo schema sulle
lezioni successive.
\end{enumerate}


\end{document}